\chapter{Manuscript-Style Research Product}

\WeekHeader{9}{Manuscript-Style Research Product}{Convert the analysis into a concise scientific report.}{Draft report with abstract, methods, results, figures, and limitations.}

\section{Why This Matters}
Writing is part of analysis. A report exposes weak logic, missing methods, unclear figures, and overclaims. Week 9 turns the student's project into a product that can become a poster, internal report, thesis chapter, manuscript section, or grant preliminary result.

\section{Report Structure}
\begin{itemize}
  \item \textbf{Title}: specific, not inflated.
  \item \textbf{Abstract}: question, approach, main result, limitation, and significance.
  \item \textbf{Introduction}: system background and why chemical traits matter.
  \item \textbf{Methods}: data source, uafR workflow, enrichment settings, filtering, visualization, and literature review.
  \item \textbf{Results}: figures and tables with restrained interpretation.
  \item \textbf{Discussion}: what the results suggest, what remains uncertain, and what should be done next.
  \item \textbf{Reproducibility statement}: how to rerun the analysis and where outputs are stored.
\end{itemize}

\section{Methods Details To Include}
\begin{itemize}
  \item uafR version.
  \item R version and operating system.
  \item Chemical query list source.
  \item \texttt{categorate()} detail level, cache setting, throttle setting, and assay detail limit.
  \item Trait matrix profile, confidence threshold, and maximum traits.
  \item Any manual filtering or grouping rules.
  \item How missing data and unmatched chemicals were handled.
\end{itemize}

\section{Guided Lab}
Run the reproducibility check:

\begin{rcode}
source("training/scripts/09_reproducibility_check.R")

check_training_project(".")
\end{rcode}

Then draft the methods paragraph from script settings and validation summaries.

\section{Independent Extension}
Students complete a 4 to 6 page report. The goal is not length. The goal is a defensible chain from question to data to evidence to interpretation.

\section{Deliverables}
\begin{tabularx}{\textwidth}{p{0.25\textwidth}YY}
\DeliverableTableHeader
Draft report & Includes all required sections and at least three figures or tables. & PDF or Word export\\
Reproducibility check & Confirms project files, scripts, results, and session information. & Log file\\
Peer review notes & Student gives and receives structured feedback. & Worksheet\\
\end{tabularx}

\section{Quality Checklist}
\begin{checklistbox}{Before Week 10}
\begin{itemize}
  \item The methods section is specific enough to rerun the analysis.
  \item Results are separated from speculation.
  \item Each figure is cited in the text.
  \item The discussion includes limitations and next steps.
\end{itemize}
\end{checklistbox}
